%%%%%%%%%%%%%%%%%%%%%%%%%%%%%%%%%%%%%%%%%%%%%%%%%
%%%%   SECTION 3a — Evaluation Setup         %%%%
%%%%   Datasets, Metrics & Protocols         %%%%
%%%%%%%%%%%%%%%%%%%%%%%%%%%%%%%%%%%%%%%%%%%%%%%%%

\section{Evaluation}
\label{sec:eval}

\subsection{Experimental Setup: Datasets, Metrics and Protocols}
\label{sec:eval_setup}

% ══════════════════════════════════════════════════════════════════════════════
% 3.1.1  PHASE 1
% ══════════════════════════════════════════════════════════════════════════════
\subsubsection{Phase 1: Retrieval Evaluation}

\paragraph{Datasets and Relevance Judgements:}
Our evaluation uses the TREC BioGen 2024 collection: 4{,}194 PubMed abstracts
and 65 biomedical topics (IDs~116--180).  Full corpus and topic details are
described in §\ref{sec:phase1}.

Relevance judgements are derived from
\texttt{biogen\_2024\_\allowbreak submissions.json} --- automated BioGen~2024 system
submissions with sentence-level citation assessments
(\texttt{evidence\_relation} field).
Evidence relation value counts:
supporting=10{,}240; neutral=1{,}412; not relevant=2{,}201;
contradicting=275; invalid citation=635.

We construct two relevance scales from this data.
\textbf{Binary qrels} (used for P@10, R@100, MAP, and MRR) assign score\,=\,1
when \texttt{evidence\_relation} equals \texttt{"supporting"} and 0~otherwise, yielding
2{,}999 relevant (topic, PMID) pairs across all 65~topics (mean 46.1 relevant
docs per topic, range 8--88; every topic has at least one relevant document).
\textbf{Graded qrels} (used for NDCG@100) map supporting$\to$5,
neutral$\to$2, and all others$\to$0, taking the maximum score per
(topic, PMID) pair.  This produces 3{,}217 graded pairs (mean 49.5 relevant
docs per topic), with 2{,}999 at score\,=\,5 and 218 at score\,=\,2.

Two caveats apply.  First, these judgements are citation-based from automated
systems, not from human assessors.  Second, the mean of 46.1 relevant documents
per topic is $3\text{--}15\times$ above standard TREC tasks, so MAP and recall
values are not directly comparable to published IR benchmarks.
Additionally, PMID~37711029 (topic~141) is absent from the filtered corpus
and is excluded with a logged warning; 11~further PMIDs appearing only in
neutral citations are also excluded.

\paragraph{Split Protocol and Cross-Validation:}
We adopt a fixed odd/even split on topic IDs (as mandated by the course
instructor): 32~train queries (odd IDs: 117, 119, \ldots, 179) and
33~test queries (even IDs: 116, 118, \ldots, 180).  The corpus is shared
across both sets --- only the query set is partitioned.

Hyperparameter selection follows a strict sequential protocol:
(1)~query field ablation using BM25 on the train set to lock the query field;
(2)~LM-JM $\lambda$ sweep on the train set with 5-fold CV;
(3)~BM25 $k_1$/$b$ grid sweep on the train set with 5-fold CV;
(4)~LM-Dir $\mu$ sweep on the train set with 5-fold CV;
(5)~dense encoder and pooling comparison on the train set via brute-force
exact cosine; and finally
(6)~a single evaluation of all locked strategies on the test set, with no
further tuning.

The 5-fold cross-validation over 32~train queries assigns roughly
26~queries to training and 6~to validation in each fold.
The primary selection criterion throughout is mean NDCG@100 across folds.

\paragraph{Metrics:}
All metrics are implemented from scratch
in \texttt{src/\allowbreak evaluation/\allowbreak metrics\allowbreak .py}
(plain Python + NumPy, no external IR library ---
following the Lab03 protocol).

\begin{itemize}
  \item \textbf{P@10} (Precision at rank~10, binary qrels) ---
        of the first 10 retrieved documents, how many are truly relevant?
        A high P@10 means the practitioner sees mostly useful references
        on the first results page without manual filtering.

  \item \textbf{R@100} (Recall at rank~100, binary qrels) ---
        of all relevant documents in the collection, how many appear in
        the top~100?  This is critical for Phase~2: any relevant study
        \emph{not} retrieved here is permanently lost to the generation
        stage.

  \item \textbf{PR Curves} --- 11-point interpolated precision--recall
        curves averaged across test queries (Lab03 pattern).  They
        visualise the precision--recall trade-off across rank thresholds;
        the area under each curve approximates Average Precision.

  \item \textbf{AP\,/\,MAP} (Average Precision / Mean Average Precision,
        binary qrels) --- $\text{AP}=\frac{1}{|\mathcal{R}|}\sum_{k=1}^{n}
        P(k)\cdot\text{rel}(k)$, averaged across queries to obtain MAP.
        AP rewards systems that rank relevant documents near the top,
        not merely retrieving them somewhere in the list.

  \item \textbf{MRR} (Mean Reciprocal Rank, binary qrels) ---
        $\text{MRR}=\frac{1}{|Q|}\sum_{q}\frac{1}{\text{rank}_q^{(1)}}$.
        How quickly does the user find the \emph{first} useful result?
        MRR\,=\,1.0 means the top-ranked document is always relevant.

  \item \textbf{NDCG@100} (Normalised Discounted Cumulative Gain, graded
        qrels) --- the \textbf{primary TREC BioGen metric}:
        $\text{DCG}=\sum_{i=1}^{100}\frac{2^{\text{rel}_i}-1}{\log_2(i+2)}$;\;
        $\text{NDCG}=\text{DCG}/\text{IDCG}$.
        Unlike binary metrics, NDCG distinguishes \emph{how} relevant a
        document is (supporting\,=\,5 vs.\ neutral\,=\,2) and penalises
        highly relevant documents placed at low ranks, making it the most
        comprehensive single number for model selection.
\end{itemize}

% ══════════════════════════════════════════════════════════════════════════════
% 3.1.2  PHASE 2
% ══════════════════════════════════════════════════════════════════════════════
\subsubsection{Phase 2: Generation Evaluation}
To be developed in future work.

% ══════════════════════════════════════════════════════════════════════════════
% 3.1.3  PHASE 3
% ══════════════════════════════════════════════════════════════════════════════
\subsubsection{Phase 3: Agentic Evaluation}
To be developed in future work.
